% Part: first-order-logic
% Chapter: beyond
% Section: other-logics

\documentclass[../../../include/open-logic-section]{subfiles}

\begin{document}

\olfileid[fr]{fol}{byd}{oth}

\olsection{Autres logiques}

Vous l'aurez peut-être compris : il n'est pas difficile d'inventer
une nouvelle logique. Vous aussi, vous pouvez définir votre propre
syntaxe, imaginer un système déductif et lui donner une sémantique.
Il vous faudra peut-être quelque ingéniosité pour que le système de
!!{derivation} soit complet pour cette sémantique, et des efforts
pour convaincre les autres que votre logique présente un véritable
intérêt. En retour, vous pourrez consacrer bien des heures
agréables à explorer ses propriétés mathématiques et informatiques.

Les dernières décennies ont vu une véritable prolifération des
logiques formelles. La logique floue vise à représenter les
raisonnements portant sur des propriétés vagues, et la logique
probabiliste ceux qui portent sur l'incertitude. Les logiques
par défaut et les logiques non monotones formalisent des raisonnements
révisables~: des inférences «~raisonnables~» peuvent être remises en cause
par de nouvelles informations. Il existe des logiques épistémiques,
consacrées aux raisonnements sur la connaissance; des logiques
causales, qui traitent des relations de cause à effet; et même
des logiques « déontiques », qui portent sur les obligations
morales et éthiques. Selon que leur motivation première est
philosophique, mathématique ou informatique, on rencontrera
l'étude de ces systèmes en logique mathématique, en logique
philosophique, en intelligence artificielle, en sciences
cognitives ou dans d'autres domaines.

La liste pourrait s'allonger encore, tant les possibilités
semblent vastes. Peut-être ne réaliserons-nous jamais le rêve
de Leibniz, qui voulait ramener l'ensemble du raisonnement
humain au calcul; cela ne nous empêche pas d'essayer.

\end{document}
